(a) Nonzero degree $d$ homogeneous polynomials, modulo multiplication by a nonzero scalar, are the points of $\mathbb{P}^N$, where $N=\binom{d+2}{2}-1$. Their zero sets give the required correspondence. By the Nullstellensatz and unique factorization, two polynomials have the same zero set exactly when they have the same distinct irreducible factors. If one is square-free of degree $d$, equality of degrees forces all the multiplicities in the other to be $1$, so the corresponding point of $\mathbb{P}^N$ is unique. Ambiguity can occur only with repeated factors.

(b) Apply elimination theory to $f,f_x,f_y,f_z$. Their having a common projective zero is a closed condition on the coefficients of $f$, since all their coefficients depend linearly on those of $f$. Its complement is nonempty by \exref{I.5.5}. By \exref{I.5.8} and \exref{I.5.9}, this complement parametrizes exactly the irreducible nonsingular degree $d$ curves, and (a) shows that the correspondence is one-to-one.
